\section{\texorpdfstring{Throughput and delay}{Throughput and delay}}\label{sec:resultados-atraso}
% FONTE: hardware/resultados_servidor/atraso_por_bloco.xlsx (fórmulas) e latencia_por_bloco.csv,
%   gerados por hardware/latencia_por_bloco.py e hardware/atraso_por_bloco_xlsx.py a partir dos
%   logs xsim do servidor (Vivado 2023.2). Os logs têm só 3 lotes (ou 20 matrizes); o intervalo
%   sai de 2 (ou 19) medidas. Faixas entre formatos: a 455 passa de T_c de +3,2 % (<18,8>) a
%   +11,5 % (<30,14>); a 152 passa de T_c por cerca de 21,5 % em todos; a folga da 606 cai de
%   11 % (<18,8>) para 3 % (<30,14>). Coleta por condição: atraso_por_bloco.xlsx (t_coleta).
%   O CSV marca o <30,14> como AP_RND_ZERO/AP_WRAP, mas o projeto é AP_TRN/AP_SAT (mesmo
%   intervalo da resultados_vivado_servidor.csv).

The architectures were designed for the throughput limit, and meeting a delay limit was outside the scope (Section~\ref{subsec:restricao-vazao}). This section shows the resulting delay and why it does not meet the latency expected for 5G. The total delay goes from the first noise sample of the block, collected at the same time in all PRBs, to the inverse of the last PRB. It has two parts: the sample collection and the time to invert the block (Equation~\ref{eq:fim-tarefas}). Table~\ref{tab:atraso-bloco} shows the two parts, from the lowest to the highest value among the formats, and Table~\ref{tab:atraso-formatos} (Appendix~\ref{ap:recursos}) shows the total delay of each format.

\begin{table}[!ht]
\caption{Throughput and delay of a block of $N_{\text{PRB}} = 275$ inversions.}
\label{tab:atraso-bloco}
\vspace{4pt}
\centering
\footnotesize
\renewcommand{\arraystretch}{1.2}
\setlength{\tabcolsep}{4pt}
\begin{tabular}{lccccc}
\toprule
\begin{tabular}[c]{@{}c@{}}\textbf{Test}\\\textbf{condition}\end{tabular} & \begin{tabular}[c]{@{}c@{}}{\boldmath $T_c$}\\\textbf{(ms)}\end{tabular} & \begin{tabular}[c]{@{}c@{}}\textbf{Block input}\\\textbf{(ms)}\end{tabular} & \begin{tabular}[c]{@{}c@{}}\textbf{Block end}\\\textbf{(ms)}\end{tabular} & \begin{tabular}[c]{@{}c@{}}\textbf{Collection}\\\textbf{(ms)}\end{tabular} & \begin{tabular}[c]{@{}c@{}}\textbf{Total delay}\\\textbf{(ms)}\end{tabular} \\
\midrule
TDLA30-5    & 25.000 & 19.473--23.724 & 19.473--23.724 & 2.000 & 21.473--25.724 \\
TDLA30-10   & 12.500 & 7.604--9.130   & 7.653--9.186   & 2.000 & 9.653--11.186 \\
TDLA30-75   & 1.667  & 1.336--1.460   & 1.485--1.615   & 0.500 & 1.985--2.115 \\
TDLC300-100 & 1.250  & 1.115--1.213   & 1.290--1.393   & 1.250 & 2.540--2.643 \\
TDLA30-300  & 0.417  & 0.369          & 0.506--0.507   & 0.417 & 0.923--0.924 \\
TDLB100-400 & 0.313  & 0.184          & 0.265--0.266   & 0.313 & 0.578--0.579 \\
\bottomrule
\end{tabular}
\vspace{4pt}
\legend{Source: the author, from the netlist simulation at 100~MHz. Block input: time to receive the 275 matrices. Block end: time until the last inversion of the block. TDLB100-400 without $\langle 30,14\rangle$, which does not fit in the FPGA.}
\end{table}

The sample collection is set by the estimation method, not by the hardware (Section~\ref{subsec:pipeline}). With 48 DM-RS REs per PRB and per slot (Section~\ref{subsubsec:disponibilidade-amostras}), it takes $\lceil N/48 \rceil$ slots, limited to $T_c$. In TDLA30-5 and TDLA30-10, the collection alone takes 2~ms, before any inversion.

The ITU-R expects that IMT-2020 can offer a latency of 1~ms over the air interface, which is the radio network's contribution to the time between sending and receiving a packet \cite[pp.~13--14]{itu-m2083}. Processing at the BS is one part of this latency \cite[Sec.~7.2.6 and Tab.~3]{itu-m2412}. Only in TDLA30-300 and TDLB100-400 is the delay below 1~ms. This delay also assumes that every slot carries the PUSCH, since the DM-RS is only transmitted together with it (Section~\ref{subsec:simbolos}). In the n78 and n257 bands, which use TDD \cite[Tabs.~5.2-1 and 5.2-2]{3gpp-ts38104}, only the uplink slots carry the PUSCH, and the collection takes longer or gathers fewer samples.

Since all architectures meet the cycle limit (Section~\ref{sec:resultados-hardware}), the whole block enters the module before $T_c$, and no queue forms. In TDLC300-100 and TDLA30-300, the last inversion of the block ends after $T_c$. This excess is a fixed delay, of about one batch latency, and it does not accumulate, because the next block enters the module on time. This holds for all formats.

The time to invert the block depends on the architecture. The format only changes this time in architectures without systolic arrays, where the interval between inversions grows with the format (Section~\ref{subsec:resultados-recursos}). The architectures of the low-mobility conditions have less parallelism and spread the 275 inversions over most of $T_c$. Thus, they use less area and power (Section~\ref{sec:resultados-hardware}), but the last PRB is ready later. More parallelism would reduce this time: the TDLB100-400 architecture, for example, finishes a block in 0.27~ms. The other PRBs get their inverses earlier, because their matrices enter the module first.

The inversion of $\mathbf G_c$, however, is only one part of the equalization, and equalization is only one of the baseband tasks of the BS (Section~\ref{subsec:uplink}). A complete delay analysis would need to add the time of the other tasks, of the covariance estimation, of the other steps of BCD-MMSE, and of the data exchange between the DUs, which this work does not consider (Section~\ref{subsec:restricao-vazao}).
